% 2021 全国甲卷文科 7：正视图与四个候选侧视图
% 绘制依据：源 PDF 第 1 页原图，400 dpi 放大逐个量取。五个同样大小的正方形。
% 关键几何：所有截线都是【边中点到边中点】——与题设 E、F、G 为过顶点 A 的三条棱
% 的中点相吻合，故截面在各投影面上的像必落在相邻两边的中点连线上。
%   正视图：【虚线】自左边中点降至底边中点，切掉【左下角】。
%   (A) 实线，自底边中点升至右边中点，切【右下角】
%   (B) 虚线，自左边中点降至底边中点，切【左下角】——与正视图同形，是「照抄正视图」的陷阱项
%   (C) 虚线，自底边中点升至右边中点，切【右下角】
%   (D) 实线，自左边中点降至底边中点，切【左下角】
% ⚠️ 初版曾把正视图与 (B)(D) 的下端画到约 0.88 宽处而非底边中点，属实质失真
%    （会改变「切哪个角」的判读）；已按 400 dpi 原图量取修正为严格中点。
% 选项标签 (A)~(D) 置于各正方形左下角外侧。
\begin{tikzpicture}[font=\small, line width=0.4pt]
  \def\s{1.75}% 正方形边长
  % 两种截线：切左下角（左边中点→底边中点）、切右下角（底边中点→右边中点）
  % #1 线型选项
  \newcommand{\cutLL}[1][]{\draw[#1] (0,{0.5*\s}) -- ({0.5*\s},0);}
  \newcommand{\cutLR}[1][]{\draw[#1] ({0.5*\s},0) -- (\s,{0.5*\s});}
  % ---------- 正视图 ----------
  \begin{scope}[shift={({1.6*\s},{1.75*\s})}]
    \draw (0,0) rectangle (\s,\s);
    \cutLL[dashed]
    \node[anchor=north] at ({0.5*\s},-0.10) {正视图};
  \end{scope}
  % ---------- 四个选项 ----------
  \begin{scope}[shift={(0,0)}]
    \draw (0,0) rectangle (\s,\s);  \cutLR
    \node[anchor=east] at (-0.06,{0.06*\s}) {(A)};
  \end{scope}
  \begin{scope}[shift={({1.55*\s},0)}]
    \draw (0,0) rectangle (\s,\s);  \cutLL[dashed]
    \node[anchor=east] at (-0.06,{0.06*\s}) {(B)};
  \end{scope}
  \begin{scope}[shift={({3.10*\s},0)}]
    \draw (0,0) rectangle (\s,\s);  \cutLR[dashed]
    \node[anchor=east] at (-0.06,{0.06*\s}) {(C)};
  \end{scope}
  \begin{scope}[shift={({4.65*\s},0)}]
    \draw (0,0) rectangle (\s,\s);  \cutLL
    \node[anchor=east] at (-0.06,{0.06*\s}) {(D)};
  \end{scope}
\end{tikzpicture}
